% !TEX program = xelatex
% Dokumen latihan matematis mandiri. Tidak di-input ke main.tex skripsi.
\documentclass[12pt,a4paper]{report}

\usepackage{amsmath,amssymb,mathtools,bm}
\usepackage{geometry}
\usepackage{enumitem}
\usepackage{hyperref}
\usepackage[provide=*,indonesian]{babel}
\geometry{left=3cm,right=3cm,top=3cm,bottom=3cm}
\hypersetup{colorlinks=true,linkcolor=blue,urlcolor=blue}
\allowdisplaybreaks

\newcommand{\dd}{\mathrm{d}}
\newcommand{\pd}{\partial}
\newcommand{\Lag}{\mathcal{L}}
\newcommand{\Riem}{R^{\rho}{}_{\sigma\mu\nu}}
\newcommand{\order}{\mathcal{O}}
\newcommand{\e}{\mathbf{e}}
% Schutz, Bab 3: basis one-form ditulis \widetilde{\omega}^{\alpha}.
\newcommand{\dual}{\widetilde{\omega}}
\newcommand{\oneformp}{\widetilde{p}}
\newcommand{\ddtilde}{\widetilde{\mathrm d}}
\newcommand{\Lie}{\mathcal{L}}
\newcommand{\Tr}{\operatorname{Tr}}

\title{Jalur Derivasi Matematis Relativitas Umum\\
\large Dari One-Form hingga Presesi Lense--Thirring\\
\normalsize Notasi Utama Mengikuti Schutz}
\author{Catatan latihan matematis untuk skripsi Gravity Probe B}
\date{Draft 0.2}

\begin{document}
\maketitle
\tableofcontents

\chapter*{Konvensi dan Peta Dependensi}
\addcontentsline{toc}{chapter}{Konvensi dan Peta Dependensi}

Koordinat ruang-waktu ditulis
\[
 x^\mu=(x^0,x^1,x^2,x^3)=(ct,r,\theta,\phi),
 \qquad \mu,\nu,\ldots\in\{0,1,2,3\}.
\]
Signature metrik yang digunakan di seluruh dokumen adalah
\[
 (-,+,+,+),\qquad u^\mu u_\mu=-c^2.
\]
Indeks Yunani mencakup ruang-waktu, indeks Latin $i,j,k$ hanya mencakup ruang.
Notasi dasar mengikuti Schutz: indeks koma menyatakan turunan parsial dan indeks
titik-koma menyatakan turunan kovarian,
\[
 V^\alpha{}_{,\beta}\equiv\pd_\beta V^\alpha,
 \qquad
 V^\alpha{}_{;\beta}\equiv\nabla_\beta V^\alpha.
\]
Schutz biasanya menggunakan satuan geometris $G=c=1$. Dokumen ini memulihkan
$G$ dan $c$ agar dimensi besaran Bumi dan GP-B dapat diperiksa secara langsung.
Pemulihan satuan ini tidak mengubah konvensi tensor Schutz.

\section*{Kamus notasi Schutz dan notasi umum}

\begin{center}
\begin{tabular}{lll}
\hline
Objek & Schutz/dokumen ini & Notasi lain yang umum\\
\hline
vektor basis & $\e_\alpha$ & $\pd_\alpha$ pada basis koordinat\\
basis one-form & $\dual^\alpha$ & $\vartheta^\alpha$ atau $\dd x^\alpha$\\
one-form & $\oneformp=p_\alpha\dual^\alpha$ & $\omega=\omega_\alpha\dd x^\alpha$\\
turunan parsial & $V^\alpha{}_{,\beta}$ & $\pd_\beta V^\alpha$\\
turunan kovarian & $V^\alpha{}_{;\beta}$ & $\nabla_\beta V^\alpha$\\
koneksi & $\Gamma^\alpha{}_{\mu\nu}$ & simbol Christoffel\\
\hline
\end{tabular}
\end{center}

Konvensi kelengkungan ditetapkan oleh
\[
 (\nabla_\mu\nabla_\nu-\nabla_\nu\nabla_\mu)V^\rho
 =R^\rho{}_{\sigma\mu\nu}V^\sigma.
\]
Dengan konvensi ini,
\[
 R^\rho{}_{\sigma\mu\nu}
 =\pd_\mu\Gamma^\rho{}_{\nu\sigma}
 -\pd_\nu\Gamma^\rho{}_{\mu\sigma}
 +\Gamma^\rho{}_{\mu\lambda}\Gamma^\lambda{}_{\nu\sigma}
 -\Gamma^\rho{}_{\nu\lambda}\Gamma^\lambda{}_{\mu\sigma}.
\]

Urutan dependensi matematis utama adalah
\[
\boxed{
\text{ruang vektor}
\to\text{ruang dual}
\to g_{\mu\nu}
\to g^{\mu\nu}
\to\Gamma^\rho{}_{\mu\nu}
\to R^\rho{}_{\sigma\mu\nu}
\to R_{\mu\nu},R
\to G_{\mu\nu}
\to\text{EFE}}
\]
dan untuk giroskop
\[
\boxed{
g_{\mu\nu}
\to x^\mu(\tau),u^\mu
\to S^\mu(\tau)
\to\bm\Omega_{\rm geo},\bm\Omega_{\rm LT}}.
\]

\chapter{Vektor, One-Form, dan Tensor}

\section{Ruang vektor dan basis}

Pada suatu titik $p$, vektor tangen $V\in T_p\mathcal{M}$ diekspansi terhadap basis
$\{\e_\mu\}$ sebagai
\[
 V=V^\mu\e_\mu.
\]
Jika koordinat berubah dari $x^\mu$ ke $x^{\mu'}$, basis dan komponen berubah
berlawanan:
\begin{align}
 \e_{\mu'}&=\frac{\pd x^\mu}{\pd x^{\mu'}}\e_\mu,\\
 V^{\mu'}&=\frac{\pd x^{\mu'}}{\pd x^\mu}V^\mu.
\end{align}
Akibatnya $V=V^{\mu'}\e_{\mu'}=V^\mu\e_\mu$ tetap merupakan objek geometris
yang sama.

Vektor basis koordinat dapat dipandang sebagai operator turunan:
\[
 \e_\mu=\pd_\mu\equiv\frac{\pd}{\pd x^\mu},
 \qquad V=V^\mu\pd_\mu.
\]
Ketika bekerja pada fungsi skalar $f$,
\[
 V[f]=V^\mu\pd_\mu f.
\]

\section{Ruang dual dan one-form}

Mengikuti Schutz Bab 3, istilah utama yang digunakan adalah \emph{one-form}.
Istilah \emph{kovektor} atau \emph{dual vector} merujuk pada objek yang sama.
Ruang dual $T_p^*\mathcal{M}$ adalah ruang semua pemetaan linear
\[
 \oneformp:T_p\mathcal{M}\longrightarrow\mathbb{R}.
\]
Linearitas berarti
\[
 \oneformp(aA+bB)=a\oneformp(A)+b\oneformp(B).
\]
Jika $A=A^\beta\e_\beta$, komponen one-form didefinisikan melalui nilainya pada
vektor basis,
\[
 p_\alpha\equiv\oneformp(\e_\alpha),
 \qquad
 \oneformp(A)=p_\alpha A^\alpha.
\]

Karena semua one-form juga membentuk ruang vektor, pilih basis one-form
$\{\dual^\alpha\}$. Kita menghendaki ekspansi Schutz (3.11),
\[
 \oneformp=p_\alpha\dual^\alpha.
\]
Terapkan kedua ruas pada $A=A^\beta\e_\beta$:
\begin{align}
 \oneformp(A)
 &=p_\alpha\dual^\alpha(A)\\
 &=p_\alpha\dual^\alpha(A^\beta\e_\beta)\\
 &=p_\alpha A^\beta\dual^\alpha(\e_\beta).
\end{align}
Agar hasil ini identik dengan $p_\alpha A^\alpha$ untuk setiap $p_\alpha$ dan
$A^\beta$, basis dual harus didefinisikan oleh Schutz (3.12),
\[
 \boxed{\dual^\alpha(\e_\beta)=\delta^\alpha{}_\beta.}
\]
Ini bukan asumsi tambahan: persamaan tersebut merupakan syarat yang membuat
koefisien $p_\alpha$ dalam ekspansi one-form sama dengan komponen yang diperoleh
dari $p_\alpha=\oneformp(\e_\alpha)$.

Untuk basis koordinat, $\e_\alpha=\pd/\pd x^\alpha$. Karena
\[
 x^\alpha{}_{,\beta}
 =\frac{\pd x^\alpha}{\pd x^\beta}
 =\delta^\alpha{}_\beta,
\]
perbandingan dengan definisi basis dual memberikan Schutz (3.20),
\[
 \boxed{\dual^\alpha=\ddtilde x^\alpha.}
\]
Tilde pada $\ddtilde x^\alpha$ menegaskan bahwa objek tersebut adalah one-form,
bukan perubahan kecil koordinat. Dalam notasi modern, tilde biasanya dihilangkan:
$\ddtilde x^\alpha\equiv\dd x^\alpha$.

Basis dual mengekstrak satu komponen vektor:
\[
 \dual^\alpha(A)
 =A^\beta\dual^\alpha(\e_\beta)
 =A^\beta\delta^\alpha{}_\beta
 =A^\alpha.
\]
Karena $\oneformp(A)=p_\alpha A^\alpha$ adalah skalar, komponen one-form
bertransformasi sebagai
\[
 p_{\alpha'}=\frac{\pd x^\alpha}{\pd x^{\alpha'}}p_\alpha.
\]

Gradien fungsi skalar secara alami merupakan one-form:
\[
 \ddtilde f=f_{,\alpha}\ddtilde x^\alpha,
 \qquad (\ddtilde f)(A)=A^\alpha f_{,\alpha}.
\]
Jadi $f_{,\alpha}=\pd_\alpha f$ membawa indeks bawah bukan karena konvensi visual, tetapi karena
ia menerima sebuah vektor dan menghasilkan skalar.

\section{Produk tensor dan tipe tensor}

Produk tensor vektor $V$ dan one-form $\oneformp$ ditentukan oleh bilinearitas:
\[
 V\otimes\oneformp
 =V^\alpha p_\beta\,\e_\alpha\otimes\dual^\beta.
\]
Tensor bertipe $(r,s)$ memiliki $r$ indeks atas dan $s$ indeks bawah:
\[
 T=T^{\mu_1\cdots\mu_r}{}_{\nu_1\cdots\nu_s}
 \e_{\mu_1}\otimes\cdots\otimes\e_{\mu_r}
 \otimes\dual^{\nu_1}\otimes\cdots\otimes\dual^{\nu_s}.
\]
Kontraksi satu indeks atas dan satu indeks bawah menurunkan orde tensor. Contoh:
\[
 T^\mu{}_\mu=\Tr(T),
 \qquad A^\mu{}_{\nu\mu}=B_\nu.
\]

\section{Metrik, panjang, dan operasi indeks}

Metrik adalah tensor simetris nondegenerat bertipe $(0,2)$:
\[
 g=g_{\mu\nu}\dual^\mu\otimes\dual^\nu,
 \qquad g_{\mu\nu}=g_{\nu\mu}.
\]
Elemen garis adalah
\[
 \dd s^2=g_{\mu\nu}\dd x^\mu\dd x^\nu.
\]
Matriks invers $g^{\mu\nu}$ didefinisikan oleh
\[
 g^{\mu\alpha}g_{\alpha\nu}=\delta^\mu{}_\nu.
\]
Metrik membentuk isomorfisme antara vektor dan one-form (kovektor):
\begin{align}
 V_\mu&=g_{\mu\nu}V^\nu,\\
 V^\mu&=g^{\mu\nu}V_\nu.
\end{align}
Contoh pada ruang-waktu Minkowski dalam koordinat Kartesian,
\[
 \eta_{\mu\nu}=\operatorname{diag}(-1,1,1,1),
 \]
maka
\[
 V_0=-V^0,\qquad V_i=V^i.
\]
Untuk tensor orde dua,
\[
 T^\mu{}_\nu=g^{\mu\alpha}T_{\alpha\nu},
 \qquad
 T^{\mu\nu}=g^{\mu\alpha}g^{\nu\beta}T_{\alpha\beta}.
\]

Determinan metrik ditulis $g=\det(g_{\mu\nu})$. Elemen volume invarian adalah
\[
 \dd^4V=\sqrt{-g}\,\dd^4x
\]
untuk signature $(-,+,+,+)$.

\section{Contoh: koordinat polar datar}

Bidang Euklides mempunyai
\[
 \dd \ell^2=\dd r^2+r^2\dd\phi^2,
 \qquad
 (g_{ij})=\begin{pmatrix}1&0\\0&r^2\end{pmatrix},
 \qquad
 (g^{ij})=\begin{pmatrix}1&0\\0&r^{-2}\end{pmatrix}.
\]
Walaupun ruangnya datar, komponen metrik bergantung pada $r$. Ini akan menghasilkan
simbol Christoffel tak nol, tetapi tensor Riemann tetap nol.

\chapter{Turunan Kovarian dan Koneksi Levi--Civita}

\section{Mengapa turunan parsial vektor tidak tensorial}

Untuk $V=V^\nu\e_\nu$,
\[
 \pd_\mu V=(\pd_\mu V^\nu)\e_\nu+V^\nu\pd_\mu\e_\nu.
\]
Basis dapat berubah dari titik ke titik. Definisikan koefisien koneksi melalui
\[
 \nabla_\mu\e_\nu=\Gamma^\rho{}_{\mu\nu}\e_\rho.
\]
Turunan kovarian komponen vektor kemudian menjadi, dalam notasi titik-koma Schutz,
\[
 \boxed{V^\nu{}_{;\mu}
 =V^\nu{}_{,\mu}+\Gamma^\nu{}_{\lambda\mu}V^\lambda}
 \quad\Longleftrightarrow\quad
 \nabla_\mu V^\nu=\pd_\mu V^\nu+\Gamma^\nu{}_{\lambda\mu}V^\lambda.
\]
Untuk menjaga aturan Leibniz pada skalar $p_\nu V^\nu$,
\begin{align}
 (p_\nu V^\nu)_{,\mu}
 &=p_{\nu;\mu}V^\nu+p_\nu V^\nu{}_{;\mu},
\end{align}
sehingga
\[
 \boxed{p_{\nu;\mu}
 =p_{\nu,\mu}-\Gamma^\lambda{}_{\nu\mu}p_\lambda}
 \quad\Longleftrightarrow\quad
 \nabla_\mu p_\nu=\pd_\mu p_\nu-\Gamma^\lambda{}_{\nu\mu}p_\lambda.
\]
Setiap indeks atas memberi suku $+\Gamma$ dan setiap indeks bawah memberi suku
$-\Gamma$. Sebagai contoh,
\[
 T^\mu{}_{\nu;\lambda}
 =T^\mu{}_{\nu,\lambda}
 +\Gamma^\mu{}_{\rho\lambda}T^\rho{}_\nu
 -\Gamma^\rho{}_{\nu\lambda}T^\mu{}_\rho.
\]

\section{Derivasi simbol Christoffel}

Koneksi Levi--Civita memenuhi dua syarat:
\begin{align}
 \nabla_\lambda g_{\mu\nu}&=0
 &&\text{(kompatibilitas metrik)},\\
 \Gamma^\rho{}_{\mu\nu}&=\Gamma^\rho{}_{\nu\mu}
 &&\text{(tanpa torsi)}.
\end{align}
Ekspansi kompatibilitas metrik memberi
\begin{align}
 \pd_\mu g_{\nu\sigma}
 &=\Gamma^\lambda{}_{\mu\nu}g_{\lambda\sigma}
 +\Gamma^\lambda{}_{\mu\sigma}g_{\nu\lambda},\tag{1}\\
 \pd_\nu g_{\mu\sigma}
 &=\Gamma^\lambda{}_{\nu\mu}g_{\lambda\sigma}
 +\Gamma^\lambda{}_{\nu\sigma}g_{\mu\lambda},\tag{2}\\
 \pd_\sigma g_{\mu\nu}
 &=\Gamma^\lambda{}_{\sigma\mu}g_{\lambda\nu}
 +\Gamma^\lambda{}_{\sigma\nu}g_{\mu\lambda}.\tag{3}
\end{align}
Menjumlahkan (1) dan (2), lalu mengurangkan (3), serta menggunakan simetri indeks
bawah koneksi,
\[
 \pd_\mu g_{\nu\sigma}+\pd_\nu g_{\mu\sigma}-\pd_\sigma g_{\mu\nu}
 =2\Gamma^\lambda{}_{\mu\nu}g_{\lambda\sigma}.
\]
Kalikan dengan $\tfrac12g^{\rho\sigma}$:
\[
 \boxed{
 \Gamma^\rho{}_{\mu\nu}
 =\frac12g^{\rho\sigma}
 (g_{\sigma\mu,\nu}+g_{\sigma\nu,\mu}-g_{\mu\nu,\sigma})}.
\]
Ini adalah bentuk yang digunakan Schutz pada Persamaan (6.32).

Simbol Christoffel bukan tensor. Di satu titik ia dapat dibuat nol dengan memilih
koordinat inersial lokal, sedangkan kelengkungan yang dibentuk darinya tidak dapat
dihilangkan jika ruang-waktu memang melengkung.

\section{Divergensi dan identitas determinan}

Dari identitas matriks $\pd_\mu\ln|g|=g^{\alpha\beta}\pd_\mu g_{\beta\alpha}$,
\[
 \Gamma^\nu{}_{\nu\mu}
 =\frac12g^{\alpha\beta}\pd_\mu g_{\alpha\beta}
 =\pd_\mu\ln\sqrt{-g}.
\]
Maka divergensi vektor dapat ditulis
\[
 \boxed{
 \nabla_\mu V^\mu
 =\frac{1}{\sqrt{-g}}\pd_\mu\left(\sqrt{-g}\,V^\mu\right)}.
\]

\section{Contoh Christoffel pada bidang polar}

Karena $\pd_r g_{\phi\phi}=2r$, komponen tak nol adalah
\begin{align}
 \Gamma^r{}_{\phi\phi}
 &=-\frac12\pd_r(r^2)=-r,\\
 \Gamma^\phi{}_{r\phi}=\Gamma^\phi{}_{\phi r}
 &=\frac12g^{\phi\phi}\pd_r g_{\phi\phi}=\frac1r.
\end{align}

\chapter{Geodesik dan Transportasi Paralel}

\section{Geodesik dari variasi panjang lintasan}

Gunakan Lagrangian kuadratik
\[
 \Lag=\frac12g_{\mu\nu}\dot x^\mu\dot x^\nu,
 \qquad \dot x^\mu=\frac{\dd x^\mu}{\dd\lambda}.
\]
Persamaan Euler--Lagrange adalah
\[
 \frac{\dd}{\dd\lambda}\left(\frac{\pd\Lag}{\pd\dot x^\rho}\right)
 -\frac{\pd\Lag}{\pd x^\rho}=0.
\]
Hitung masing-masing suku:
\begin{align}
 \frac{\pd\Lag}{\pd\dot x^\rho}&=g_{\rho\nu}\dot x^\nu,\\
 \frac{\dd}{\dd\lambda}(g_{\rho\nu}\dot x^\nu)
 &=g_{\rho\nu}\ddot x^\nu
 +(\pd_\sigma g_{\rho\nu})\dot x^\sigma\dot x^\nu,\\
 \frac{\pd\Lag}{\pd x^\rho}
 &=\frac12(\pd_\rho g_{\mu\nu})\dot x^\mu\dot x^\nu.
\end{align}
Kalikan hasilnya dengan $g^{\alpha\rho}$ dan simetriskan pasangan
$\dot x^\mu\dot x^\nu$:
\[
 \boxed{
 \ddot x^\alpha+\Gamma^\alpha{}_{\mu\nu}\dot x^\mu\dot x^\nu=0}.
\]
Parameter $\lambda$ yang membuat ruas kanan nol disebut parameter afin. Untuk
partikel bermassa dapat dipilih $\lambda=\tau$. Jika memakai parameter nonafin,
\[
 \ddot x^\alpha+\Gamma^\alpha{}_{\mu\nu}\dot x^\mu\dot x^\nu
 =f(\lambda)\dot x^\alpha.
\]

\section{Transportasi paralel}

Sepanjang kurva $x^\mu(\lambda)$ dengan vektor singgung
$u^\mu=\dd x^\mu/\dd\lambda$, turunan kovarian sepanjang kurva adalah
\[
 \frac{D}{D\lambda}=u^\mu\nabla_\mu.
\]
Vektor $V^\mu$ tertransportasi paralel jika
\[
 \boxed{
 \frac{DV^\mu}{D\lambda}
 =\frac{\dd V^\mu}{\dd\lambda}
 +\Gamma^\mu{}_{\alpha\beta}u^\alpha V^\beta=0}.
\]
Geodesik adalah kasus khusus ketika vektor yang ditransportasikan adalah vektor
singgung itu sendiri:
\[
 \frac{Du^\mu}{D\lambda}=0.
\]

Kompatibilitas metrik menjamin produk dalam dipertahankan. Jika $V$ dan $W$
keduanya ditransportasikan paralel,
\begin{align}
 \frac{\dd}{\dd\lambda}(g_{\mu\nu}V^\mu W^\nu)
 &=u^\alpha\nabla_\alpha(g_{\mu\nu}V^\mu W^\nu)=0.
\end{align}

Untuk giroskop ideal yang bergerak geodesik,
\begin{align}
 \frac{DS^\mu}{D\tau}&=0,\\
 \frac{Du^\mu}{D\tau}&=0,\\
 S_\mu u^\mu&=0.
\end{align}
Kekekalan ortogonalitas mengikuti langsung:
\[
 \frac{D}{D\tau}(S_\mu u^\mu)
 =\frac{DS_\mu}{D\tau}u^\mu+S_\mu\frac{Du^\mu}{D\tau}=0.
\]
Jika lintasan memiliki empat-percepatan $a^\mu=Du^\mu/D\tau\ne0$, transportasi
spin tanpa rotasi intrinsik diberikan oleh transportasi Fermi--Walker:
\[
 \frac{DS^\mu}{D\tau}
 =\frac{1}{c^2}(u^\mu a_\nu-a^\mu u_\nu)S^\nu.
\]
Ruas kanan kembali nol pada lintasan geodesik.

\section{Holonomi}

Transportasi paralel mengelilingi kurva tertutup umumnya tidak mengembalikan vektor
ke orientasi semula. Untuk loop infinitesimal dengan elemen luas bivector
$\Sigma^{\mu\nu}$,
\[
 \delta V^\rho
 =-\frac12R^\rho{}_{\sigma\mu\nu}V^\sigma\Sigma^{\mu\nu}
 +\order(\Sigma^{3/2}).
\]
Jadi holonomi bukan objek tambahan yang terpisah: ia merupakan aksi terintegrasi
dari kelengkungan pada vektor yang ditransportasikan.

\chapter{Tensor Riemann, Ricci, dan Geodesic Deviation}

\section{Riemann dari komutator turunan kovarian}

Untuk vektor $V^\alpha$,
\begin{align}
 \nabla_\mu\nabla_\nu V^\rho
 &=\pd_\mu(\pd_\nu V^\rho+\Gamma^\rho{}_{\nu\sigma}V^\sigma)
 +\Gamma^\rho{}_{\mu\lambda}(\pd_\nu V^\lambda
 +\Gamma^\lambda{}_{\nu\sigma}V^\sigma)
 -\Gamma^\lambda{}_{\mu\nu}(\pd_\lambda V^\rho
 +\Gamma^\rho{}_{\lambda\sigma}V^\sigma).
\end{align}
Kurangkan ekspresi dengan $\mu\leftrightarrow\nu$. Suku yang memuat turunan
$V^\rho$ saling hapus untuk koneksi tanpa torsi, sehingga
\[
 (\nabla_\mu\nabla_\nu-\nabla_\nu\nabla_\mu)V^\alpha
 =R^\alpha{}_{\beta\mu\nu}V^\beta,
\]
dengan bentuk Schutz (6.63)
\[
 \boxed{
 R^\alpha{}_{\beta\mu\nu}
 =\Gamma^\alpha{}_{\beta\nu,\mu}
 -\Gamma^\alpha{}_{\beta\mu,\nu}
 +\Gamma^\alpha{}_{\sigma\mu}\Gamma^\sigma{}_{\beta\nu}
 -\Gamma^\alpha{}_{\sigma\nu}\Gamma^\sigma{}_{\beta\mu}}.
\]

Turunkan indeks pertama dengan metrik:
\[
 R_{\rho\sigma\mu\nu}=g_{\rho\alpha}R^\alpha{}_{\sigma\mu\nu}.
\]
Simetri-simetrinya adalah
\begin{align}
 R_{\rho\sigma\mu\nu}&=-R_{\sigma\rho\mu\nu},\\
 R_{\rho\sigma\mu\nu}&=-R_{\rho\sigma\nu\mu},\\
 R_{\rho\sigma\mu\nu}&=R_{\mu\nu\rho\sigma},\\
 R_{\rho[\sigma\mu\nu]}&=0.
\end{align}

\section{Kontraksi Ricci, skalar Ricci, dan Einstein}

Tensor Ricci diperoleh dengan mengontraksi indeks pertama dan ketiga, sesuai
Schutz (6.91):
\[
 \boxed{R_{\alpha\beta}=R^\mu{}_{\alpha\mu\beta}}.
\]
Dalam simbol Christoffel,
\[
 R_{\mu\nu}
 =\pd_\rho\Gamma^\rho{}_{\nu\mu}
 -\pd_\nu\Gamma^\rho{}_{\rho\mu}
 +\Gamma^\rho{}_{\rho\lambda}\Gamma^\lambda{}_{\nu\mu}
 -\Gamma^\rho{}_{\nu\lambda}\Gamma^\lambda{}_{\rho\mu}.
\]
Skalar Ricci dan tensor Einstein adalah
\begin{align}
 R&=g^{\mu\nu}R_{\mu\nu},\\
 \boxed{G_{\mu\nu}=R_{\mu\nu}-\frac12Rg_{\mu\nu}}.
\end{align}

Identitas Bianchi diferensial
\[
 \nabla_{[\lambda}R_{\rho\sigma]\mu\nu}=0
\]
setelah kontraksi dua kali menghasilkan
\[
 \boxed{\nabla_\mu G^{\mu\nu}=0}.
\]
Inilah alasan geometris mengapa ruas materi pada persamaan Einstein harus memenuhi
$\nabla_\mu T^{\mu\nu}=0$.

\section{Geodesic deviation}

Pertimbangkan keluarga geodesik $x^\mu(\tau,s)$. Definisikan
\[
 u^\mu=\frac{\pd x^\mu}{\pd\tau},
 \qquad \xi^\mu=\frac{\pd x^\mu}{\pd s},
 \qquad [u,\xi]=0.
\]
Karena $\nabla_u u=0$ dan $\nabla_u\xi=\nabla_\xi u$,
\begin{align}
 \nabla_u\nabla_u\xi
 &=\nabla_u\nabla_\xi u\\
 &=\nabla_\xi\nabla_u u+R(u,\xi)u\\
 &=R(u,\xi)u.
\end{align}
Dalam komponen,
\[
 \boxed{
 \frac{D^2\xi^\mu}{D\tau^2}
 =R^\mu{}_{\nu\alpha\beta}u^\nu u^\alpha\xi^\beta}.
\]
Tanda ruas kanan mengikuti konvensi Riemann yang ditetapkan pada awal dokumen.
Persamaan ini menghubungkan kelengkungan dengan percepatan relatif dua partikel
jatu bebas yang berdekatan.

\section{Cek: bidang polar tetap datar}

Gunakan $\Gamma^r{}_{\phi\phi}=-r$ dan
$\Gamma^\phi{}_{r\phi}=1/r$. Salah satu komponen potensial adalah
\begin{align}
 R^r{}_{\phi r\phi}
 &=\pd_r\Gamma^r{}_{\phi\phi}
 -\pd_\phi\Gamma^r{}_{r\phi}
 +\Gamma^r{}_{r\lambda}\Gamma^\lambda{}_{\phi\phi}
 -\Gamma^r{}_{\phi\lambda}\Gamma^\lambda{}_{r\phi}\\
 &=-1-(-r)(1/r)=0.
\end{align}
Christoffel tak nol tidak cukup untuk menyimpulkan adanya kelengkungan.

\chapter{Persamaan Medan Einstein}

\section{Aksi Einstein--Hilbert}

Aksi gravitasi dan materi adalah
\[
 S=S_g+S_m
 =\frac{c^3}{16\pi G}\int(R-2\Lambda)\sqrt{-g}\,\dd^4x+S_m.
\]
Variasi dilakukan terhadap $g^{\mu\nu}$. Gunakan
\begin{align}
 \delta\sqrt{-g}&=-\frac12\sqrt{-g}\,g_{\mu\nu}\delta g^{\mu\nu},\\
 \delta R&=R_{\mu\nu}\delta g^{\mu\nu}
 +g^{\mu\nu}\delta R_{\mu\nu}.
\end{align}
Identitas Palatini,
\[
 \delta R_{\mu\nu}
 =\nabla_\rho\delta\Gamma^\rho{}_{\nu\mu}
 -\nabla_\nu\delta\Gamma^\rho{}_{\rho\mu},
\]
membuat $g^{\mu\nu}\delta R_{\mu\nu}$ menjadi suku batas. Setelah suku batas
dihilangkan atau ditangani dengan suku Gibbons--Hawking--York,
\[
 \delta S_g
 =\frac{c^3}{16\pi G}\int
 (G_{\mu\nu}+\Lambda g_{\mu\nu})
 \delta g^{\mu\nu}\sqrt{-g}\,\dd^4x.
\]
Definisikan tensor energi--momentum melalui
\[
 \delta S_m=-\frac{1}{2c}\int T_{\mu\nu}\delta g^{\mu\nu}
 \sqrt{-g}\,\dd^4x.
\]
Syarat $\delta S=0$ untuk variasi bebas menghasilkan
\[
 \boxed{
 G_{\mu\nu}+\Lambda g_{\mu\nu}
 =\frac{8\pi G}{c^4}T_{\mu\nu}}.
\]

\section{Trace-reversed form dan vakum}

Ambil trace dalam empat dimensi:
\begin{align}
 g^{\mu\nu}G_{\mu\nu}+4\Lambda
 &=\frac{8\pi G}{c^4}T,\\
 -R+4\Lambda&=\frac{8\pi G}{c^4}T.
\end{align}
Substitusi kembali memberi
\[
 R_{\mu\nu}
 =\frac{8\pi G}{c^4}\left(T_{\mu\nu}-\frac12Tg_{\mu\nu}\right)
 +\Lambda g_{\mu\nu}.
\]
Untuk daerah vakum tanpa konstanta kosmologis,
\[
 T_{\mu\nu}=0,\quad\Lambda=0
 \quad\Longrightarrow\quad
 \boxed{R_{\mu\nu}=0}.
\]
Vakum tidak mengharuskan $R^\rho{}_{\sigma\mu\nu}=0$; Schwarzschild dan Kerr
adalah contoh $R_{\mu\nu}=0$ tetapi Riemann tidak nol.

\section{Limit Newtonian}

Ambil gangguan lemah dan statik
\[
 g_{\mu\nu}=\eta_{\mu\nu}+h_{\mu\nu},
 \qquad |h_{\mu\nu}|\ll1,
 \qquad \pd_0h_{\mu\nu}\approx0.
\]
Untuk gerak lambat, geodesik spasial didominasi oleh
\begin{align}
 \frac{\dd^2x^i}{\dd t^2}
 &\simeq-c^2\Gamma^i{}_{00}
 =\frac{c^2}{2}\pd_i h_{00}.
\end{align}
Bandingkan dengan $\ddot{\bm x}=-\bm\nabla\Phi$:
\[
 \boxed{h_{00}=-\frac{2\Phi}{c^2}},
 \qquad
 g_{00}=-\left(1+\frac{2\Phi}{c^2}\right).
\]
Untuk materi nonrelativistik, $T_{00}\simeq\rho c^2$ dan EFE linear menghasilkan
\[
 \nabla^2\Phi=4\pi G\rho.
\]

\chapter{Derivasi Metrik Schwarzschild}

\section{Ansatz statik dan simetri bola}

Elemen garis paling umum setelah memilih koordinat areal $r$ dapat ditulis
\[
 \dd s^2=-A(r)c^2\dd t^2+B(r)\dd r^2
 +r^2(\dd\theta^2+\sin^2\theta\,\dd\phi^2).
\]
Komponen inversnya adalah
\[
 g^{tt}=-\frac{1}{A c^2},\quad
 g^{rr}=\frac1B,\quad
 g^{\theta\theta}=\frac1{r^2},\quad
 g^{\phi\phi}=\frac1{r^2\sin^2\theta}.
\]

\section{Christoffel yang diperlukan}

Tanda prima berarti $\dd/\dd r$. Komponen independen tak nol adalah
\begin{align}
 \Gamma^t{}_{tr}&=\Gamma^t{}_{rt}=\frac{A'}{2A},\\
 \Gamma^r{}_{tt}&=\frac{c^2A'}{2B},
 &\Gamma^r{}_{rr}&=\frac{B'}{2B},\\
 \Gamma^r{}_{\theta\theta}&=-\frac rB,
 &\Gamma^r{}_{\phi\phi}&=-\frac{r\sin^2\theta}{B},\\
 \Gamma^\theta{}_{r\theta}&=\Gamma^\phi{}_{r\phi}=\frac1r,
 &\Gamma^\theta{}_{\phi\phi}&=-\sin\theta\cos\theta,\\
 \Gamma^\phi{}_{\theta\phi}&=\cot\theta.
\end{align}

Substitusi ke definisi Ricci memberi
\begin{align}
 R_{tt}&=\frac{c^2}{2B}\left[
 A''-\frac{A'B'}{2B}-\frac{(A')^2}{2A}+\frac{2A'}r
 \right],\\
 R_{rr}&=-\frac{A''}{2A}+\frac{(A')^2}{4A^2}
 +\frac{A'B'}{4AB}+\frac{B'}{Br},\\
 R_{\theta\theta}&=1-\frac1B
 +\frac{rB'}{2B^2}-\frac{rA'}{2AB},\\
 R_{\phi\phi}&=\sin^2\theta\,R_{\theta\theta}.
\end{align}

\section{Penyelesaian persamaan vakum}

Dari kombinasi $R_{tt}/(Ac^2)+R_{rr}/B=0$ diperoleh
\[
 \frac{A'}{A}+\frac{B'}B=0
 \quad\Longrightarrow\quad AB=C.
\]
Asimtotik datar mensyaratkan $A\to1$ dan $B\to1$ saat $r\to\infty$, sehingga
$C=1$ dan
\[
 B=A^{-1}.
\]
Masukkan ke $R_{\theta\theta}=0$:
\begin{align}
 0&=1-A-rA',\\
 \frac{\dd}{\dd r}(rA)&=1,\\
 rA&=r+C_1,\\
 A(r)&=1+\frac{C_1}{r}.
\end{align}
Limit Newtonian memakai $\Phi=-GM/r$ dan
$g_{tt}=-A c^2\simeq-(1+2\Phi/c^2)c^2$, maka
\[
 C_1=-\frac{2GM}{c^2}\equiv-r_s.
\]
Dengan $r_s=2GM/c^2$,
\[
 \boxed{
 \dd s^2=-\left(1-\frac{r_s}{r}\right)c^2\dd t^2
 +\left(1-\frac{r_s}{r}\right)^{-1}\dd r^2
 +r^2\dd\Omega^2}.
\]

\section{Konstanta gerak dan potensial efektif}

Karena simetri bola, orbit dapat dipilih pada $\theta=\pi/2$. Lagrangian per
satuan massa adalah
\[
 2\Lag=-f c^2\dot t^2+f^{-1}\dot r^2+r^2\dot\phi^2,
 \qquad f=1-\frac{r_s}{r}.
\]
Koordinat $t$ dan $\phi$ siklik, sehingga
\begin{align}
 E&=fc^2\dot t,\\
 L&=r^2\dot\phi.
\end{align}
Normalisasi $u^\mu u_\mu=-c^2$ memberi
\begin{align}
 -c^2&=-fc^2\dot t^2+f^{-1}\dot r^2+r^2\dot\phi^2,\\
 \dot r^2&=\frac{E^2}{c^2}-f\left(c^2+\frac{L^2}{r^2}\right).
\end{align}
Definisikan
\[
 V_{\rm eff}(r)=f\left(c^2+\frac{L^2}{r^2}\right).
\]
Orbit lingkaran memenuhi
\[
 \dot r=0,\qquad \frac{\dd V_{\rm eff}}{\dd r}=0,
\]
dan kestabilannya ditentukan oleh
\[
 \frac{\dd^2V_{\rm eff}}{\dd r^2}>0.
\]
Batas orbit lingkaran stabil terdalam adalah
\[
 r_{\rm ISCO}=\frac{6GM}{c^2}=3r_s.
\]
Untuk orbit GP-B, $r\gg r_s$, sehingga orbit berada sangat jauh dari rezim ISCO.

\section{Kretschmann scalar}

Karena $R_{\mu\nu}=0$ dan $R=0$ di vakum, kelengkungan Schwarzschild terdeteksi
oleh invarian Riemann
\[
 K=R_{\alpha\beta\gamma\delta}R^{\alpha\beta\gamma\delta}
 =\frac{48G^2M^2}{c^4r^6}.
\]
$K$ hingga pada $r=r_s$ tetapi divergen pada $r=0$, sehingga $r=r_s$ adalah
singularitas koordinat sedangkan $r=0$ adalah singularitas kelengkungan.

\chapter{Geometri Kerr dan Geodesiknya}

\section{Parameter dan fungsi dasar}

Definisikan parameter panjang
\[
 m=\frac{GM}{c^2},\qquad a=\frac{J}{Mc},
\]
serta
\[
 \rho^2=r^2+a^2\cos^2\theta,
 \qquad
 \Delta=r^2-2mr+a^2.
\]
Schutz menulis parameter massa geometris ini sebagai $M$ dan mendefinisikan
$a=J/M$ karena menggunakan $G=c=1$. Dengan satuan dipulihkan, $M_{\rm Schutz}=m$
dan $a=J/(Mc)$. Simbol $\rho^2$ dan $\Delta$ mengikuti Schutz (11.72).
Dalam koordinat Boyer--Lindquist $x^0=ct$, metrik Kerr dengan signature
$(-,+,+,+)$ adalah
\begin{align}
 \dd s^2={}&-\left(1-\frac{2mr}{\rho^2}\right)(\dd x^0)^2
 -\frac{4mar\sin^2\theta}{\rho^2}\,\dd x^0\dd\phi
 +\frac{\rho^2}{\Delta}\dd r^2+\rho^2\dd\theta^2\\
 &+\left(r^2+a^2+\frac{2ma^2r\sin^2\theta}{\rho^2}\right)
 \sin^2\theta\,\dd\phi^2.
\end{align}
Karena suku silang dalam elemen garis adalah $2g_{0\phi}\dd x^0\dd\phi$,
\[
 g_{0\phi}=-\frac{2mar\sin^2\theta}{\rho^2}.
\]
Matriks blok $(0,\phi)$ tidak diagonal. Determinan keseluruhan adalah
\[
 g=-\rho^4\sin^2\theta.
\]

\section{Invers metrik}

Untuk blok
\[
 \begin{pmatrix}g_{00}&g_{0\phi}\\g_{0\phi}&g_{\phi\phi}\end{pmatrix},
\]
invers diperoleh dengan rumus invers matriks $2\times2$. Hasilnya
\begin{align}
 g^{00}&=-\frac{(r^2+a^2)^2-a^2\Delta\sin^2\theta}
 {\rho^2\Delta},\\
 g^{0\phi}&=-\frac{2mar}{\rho^2\Delta},\\
 g^{\phi\phi}&=\frac{\Delta-a^2\sin^2\theta}
 {\rho^2\Delta\sin^2\theta},\\
 g^{rr}&=\frac{\Delta}{\rho^2},
 &g^{\theta\theta}&=\frac1{\rho^2}.
\end{align}
Pemeriksaan wajib adalah
\[
 g^{\mu\alpha}g_{\alpha\nu}=\delta^\mu{}_\nu.
\]

\section{Limit matematis}

Untuk $a\to0$,
\[
 \rho^2\to r^2,\qquad\Delta\to r(r-2m),\qquad g_{0\phi}\to0,
\]
dan metrik mereduksi tepat ke Schwarzschild. Untuk $r\to\infty$,
\[
 g_{\mu\nu}\to\eta_{\mu\nu}
\]
dalam koordinat bola. Untuk $M\to0$ tetapi $a\ne0$, metrik merupakan ruang-waktu
datar dalam koordinat sferoidal pepat.

\section{Horizon dan ergosurface}

Horizon koordinat Boyer--Lindquist ditentukan oleh $g^{rr}=0$:
\[
 \Delta=0
 \quad\Longrightarrow\quad
 r_\pm=m\pm\sqrt{m^2-a^2}.
\]
Permukaan batas statik ditentukan oleh $g_{00}=0$:
\[
 r_{\rm ergo}(\theta)=m+\sqrt{m^2-a^2\cos^2\theta}.
\]
Untuk Bumi, objek tidak berada dekat horizon; ekspresi ini dipakai sebagai
pengecekan struktur solusi, bukan rezim simulasi GP-B.

\section{Konstanta gerak}

Karena metrik tidak bergantung eksplisit pada $x^0$ dan $\phi$, terdapat dua
vektor Killing
\[
 \xi_{(t)}^\mu=(1,0,0,0),
 \qquad
 \xi_{(\phi)}^\mu=(0,0,0,1).
\]
Untuk geodesik dengan momentum $p_\mu$, kuantitas
\[
 E=-p_0c,
 \qquad L_z=p_\phi
\]
kekal. Konservasi mengikuti
\begin{align}
 \frac{\dd}{\dd\tau}(\xi_\mu u^\mu)
 &=u^\nu\nabla_\nu(\xi_\mu u^\mu)\\
 &=u^\mu u^\nu\nabla_\nu\xi_\mu
 +\xi_\mu u^\nu\nabla_\nu u^\mu\\
 &=u^\mu u^\nu\nabla_{(\nu}\xi_{\mu)}=0.
\end{align}

Untuk pemisahan berikut saja digunakan satuan geometris $G=c=1$ dan massa uji
dinormalisasi menjadi satu. Persamaan Hamilton--Jacobi
\[
 g^{\mu\nu}\pd_\mu S\,\pd_\nu S=-1
\]
dapat dipisahkan dengan ansatz
\[
 S=-Et+L_z\phi+S_r(r)+S_\theta(\theta)+\frac12\lambda.
\]
Konstanta pemisahan menghasilkan konstanta Carter $Q$. Geodesik timelike memenuhi
\begin{align}
 \rho^4\dot r^2
 &=\left[E(r^2+a^2)-aL_z\right]^2
 -\Delta\left[r^2+(L_z-aE)^2+Q\right],\\
 \rho^4\dot\theta^2
 &=Q-\cos^2\theta\left[a^2(1-E^2)
 +\frac{L_z^2}{\sin^2\theta}\right].
\end{align}
Kedua persamaan ini memperlihatkan bahwa gerak Kerr terpisah menjadi fungsi radial
dan polar, meskipun metrik memuat pencampuran koordinat waktu dan azimut.

\chapter{Medan Lemah, Gravitomagnetisme, dan Lense--Thirring}

\section{Linearisasi persamaan Einstein}

Tuliskan
\[
 g_{\mu\nu}=\eta_{\mu\nu}+h_{\mu\nu},
 \qquad |h_{\mu\nu}|\ll1.
\]
Naik-turunkan indeks $h_{\mu\nu}$ dengan $\eta_{\mu\nu}$ pada orde pertama.
Definisikan trace-reverse
\[
 \bar h_{\mu\nu}=h_{\mu\nu}-\frac12\eta_{\mu\nu}h,
 \qquad h=\eta^{\alpha\beta}h_{\alpha\beta}.
\]
Dalam gauge Lorenz
\[
 \pd^\mu\bar h_{\mu\nu}=0,
\]
EFE linear menjadi
\[
 \boxed{
 \Box\bar h_{\mu\nu}=-\frac{16\pi G}{c^4}T_{\mu\nu}},
 \qquad \Box=\eta^{\alpha\beta}\pd_\alpha\pd_\beta.
\]
Untuk sumber stasioner, $\Box\to\nabla^2$, sehingga
\[
 \bar h_{\mu\nu}(\bm x)
 =\frac{4G}{c^4}\int
 \frac{T_{\mu\nu}(\bm x')}{|\bm x-\bm x'|}\,\dd^3x'.
\]

\section{Suku massa dan arus massa}

Untuk materi lambat,
\[
 T_{00}\simeq\rho c^2,
 \qquad T_{0i}\simeq-\rho c v_i.
\]
Bagian $00$ menghasilkan potensial Newton $\Phi$. Bagian $0i$ menghasilkan potensial
vektor gravitomagnetik. Definisikan konvensi
\[
 \bm A_g(\bm x)
 =G\int\frac{\rho(\bm x')\bm v(\bm x')}
 {|\bm x-\bm x'|}\,\dd^3x'.
\]
Jauh dari benda berotasi dengan momentum sudut
\[
 \bm J=\int\bm x'\times\rho(\bm x')\bm v(\bm x')\,\dd^3x',
\]
ekspansi multipol memberi
\[
 \boxed{
 \bm A_g=\frac{G}{2}\frac{\bm J\times\bm r}{r^3}}
\]
dalam normalisasi yang dipakai di bagian ini. Medan dipol yang bersesuaian adalah
\begin{align}
 \bm B_g&=\bm\nabla\times\bm A_g\\
 &=\frac{G}{2r^3}\left[3\hat{\bm r}
 (\hat{\bm r}\cdot\bm J)-\bm J\right].
\end{align}
Faktor numerik antara $\bm A_g$, $h_{0i}$, dan $\bm B_g$ berbeda pada beberapa
konvensi gravitomagnetisme. Kuantitas akhir $\bm\Omega_{\rm LT}$ di bawah tidak
bergantung pada penamaan medan antara tersebut.

\section{Ekspansi Kerr dan asal $g_{0\phi}$}

Untuk $r\gg m$ dan $r\gg a$,
\[
 \rho^2=r^2+\order(a^2),
\]
sehingga
\begin{align}
 g_{00}&=-1+\frac{2GM}{c^2r}+\order(r^{-2},a^2),\\
 g_{0\phi}&=-\frac{2GJ}{c^3r}\sin^2\theta
 +\order(r^{-2},a^2).
\end{align}
Karena $x^0=ct$, elemen garis mengandung
\[
 2g_{0\phi}\dd x^0\dd\phi
 =-\frac{4GJ}{c^2r}\sin^2\theta\,\dd t\dd\phi.
\]
Suku ini lenyap ketika $J=0$ dan merupakan jejak matematis rotasi sumber pada
geometri eksternal.

\section{Presesi spin pada medan stasioner lemah}

Ambil tetrad lokal ortonormal pengamat dan tulis komponen spasial spin sebagai
$\bm S$. Reduksi orde pertama persamaan transportasi
\[
 \frac{\dd S^\mu}{\dd\tau}
 =-\Gamma^\mu{}_{\alpha\beta}u^\alpha S^\beta
\]
memberi bentuk rotasi
\[
 \frac{\dd\bm S}{\dd t}=\bm\Omega\times\bm S.
\]
Kontribusi yang linear terhadap momentum sudut sumber adalah
\[
 \boxed{
 \bm\Omega_{\rm LT}
 =\frac{G}{c^2r^3}\left[
 3\hat{\bm r}(\hat{\bm r}\cdot\bm J)-\bm J
 \right]}.
\]
Untuk titik pada bidang ekuator, $\hat{\bm r}\cdot\bm J=0$, sehingga
\[
 \bm\Omega_{\rm LT}=-\frac{G\bm J}{c^2r^3}.
\]
Pada sumbu rotasi, $\hat{\bm r}\parallel\bm J$, sehingga
\[
 \bm\Omega_{\rm LT}=\frac{2G\bm J}{c^2r^3}.
\]

\section{Presesi geodetik}

Bagian medan lemah yang bergantung pada massa, bersama gerak orbit giroskop,
menghasilkan
\[
 \boxed{
 \bm\Omega_{\rm geo}
 =\frac{3GM}{2c^2r^3}(\bm r\times\bm v)}.
\]
Untuk orbit lingkaran, $\bm r\perp\bm v$ dan $v=\sqrt{GM/r}$, maka
\[
 \Omega_{\rm geo}
 =\frac{3GM}{2c^2r}n,
 \qquad n=\sqrt{\frac{GM}{r^3}}.
\]
Presesi per orbit adalah
\[
 \Delta\psi_{\rm geo}
 =\Omega_{\rm geo}\frac{2\pi}{n}
 =\frac{3\pi GM}{c^2r}.
\]
Untuk orbit elips dengan semi-major axis $a_o$ dan eksentrisitas $e$,
\[
 \Delta\psi_{\rm geo}
 =\frac{3\pi GM}{c^2a_o(1-e^2)}.
\]

\section{Rata-rata orbit polar dan observabel GP-B}

Untuk orbit lingkaran polar, gunakan
\[
 \langle\hat r_i\hat r_j\rangle
 =\frac12(\delta_{ij}-\hat L_i\hat L_j),
\]
dengan $\hat{\bm L}$ normal bidang orbit. Maka
\begin{align}
 \left\langle\bm\Omega_{\rm LT}\right\rangle
 &=\frac{G}{c^2r^3}\left[
 \frac32(\bm J-(\bm J\cdot\hat{\bm L})\hat{\bm L})-\bm J
 \right].
\end{align}
Untuk orbit polar ideal, $\bm J\cdot\hat{\bm L}=0$, sehingga
\[
 \boxed{
 \left\langle\bm\Omega_{\rm LT}\right\rangle
 =\frac{G\bm J}{2c^2r^3}}.
\]
Prediksi yang dibandingkan dengan GP-B diperoleh dengan memproyeksikan vektor
presesi total ke dua sumbu pengamatan:
\[
 \bm\Omega=\bm\Omega_{\rm geo}+\bm\Omega_{\rm LT},
 \qquad
 \Omega_A=\bm\Omega\cdot\hat{\bm e}_A.
\]

\chapter{Sistem ODE Orbit dan Spin}

\section{Reduksi ke orde satu}

Definisikan state
\[
 \bm y=(x^\mu,u^\mu,S^\mu)\in\mathbb{R}^{12}.
\]
Persamaan geodesik dan transportasi paralel menjadi
\begin{align}
 \frac{\dd x^\mu}{\dd\tau}&=u^\mu,\\
 \frac{\dd u^\mu}{\dd\tau}
 &=-\Gamma^\mu{}_{\alpha\beta}(x)u^\alpha u^\beta,\\
 \frac{\dd S^\mu}{\dd\tau}
 &=-\Gamma^\mu{}_{\alpha\beta}(x)u^\alpha S^\beta.
\end{align}
Dengan
\[
 \left[A(\tau)\right]^\mu{}_{\beta}
 =\Gamma^\mu{}_{\alpha\beta}(x(\tau))u^\alpha(\tau),
\]
persamaan spin berbentuk sistem linear dengan koefisien berubah:
\[
 \dot{\bm S}=-A(\tau)\bm S.
\]
Solusi formalnya adalah exponential berurutan lintasan
\[
 \bm S(\tau)
 =\mathcal{P}\exp\left[-\int_{\tau_0}^{\tau}A(\tau')\dd\tau'\right]
 \bm S(\tau_0).
\]

\section{Constraint yang harus dipantau}

Tiga residual numerik utama adalah
\begin{align}
 C_u&=g_{\mu\nu}u^\mu u^\nu+c^2,\\
 C_S&=g_{\mu\nu}S^\mu S^\nu-S_0^2,\\
 C_\perp&=g_{\mu\nu}S^\mu u^\nu.
\end{align}
Solusi kontinu ideal memenuhi
\[
 C_u=C_S=C_\perp=0.
\]
Residual tak nol pada integrasi numerik mengukur error solver, implementasi koneksi,
atau kondisi awal yang tidak konsisten.

\section{Ekstraksi sudut presesi}

Komponen koordinat $S^i$ tidak langsung merupakan komponen yang diukur. Gunakan
tetrad ortonormal $e_{\hat a}{}^\mu$ milik pengamat:
\[
 S^{\hat a}=e^{\hat a}{}_{\mu}S^\mu,
 \qquad
 g_{\mu\nu}e_{\hat a}{}^\mu e_{\hat b}{}^\nu
 =\eta_{\hat a\hat b}.
\]
Setelah memperoleh vektor ruang lokal $\bm S_{\rm loc}$, sudut terhadap arah awal
dapat dihitung secara stabil dengan
\[
 \Delta\psi
 =\operatorname{atan2}\left(
 \|\bm S_0\times\bm S\|,
 \bm S_0\cdot\bm S
 \right).
\]
Untuk sudut kecil, $\Delta\psi\simeq\|\bm S_0\times\bm S\|/S_0^2$.

\chapter{Latihan Kepekaan Matematis}

Kerjakan tanpa melihat hasil akhir. Tujuan latihan adalah membangun kepekaan indeks,
tanda, dimensi, limit, dan dependensi antarpersamaan.

\begin{enumerate}[label=\arabic*.]
 \item Buktikan bahwa $\oneformp(V)=p_\mu V^\mu$ invarian di bawah perubahan koordinat.
 \item Turunkan aturan transformasi tensor $(1,1)$ dari transformasi basis dan basis dual.
 \item Dari $g^{\mu\alpha}g_{\alpha\nu}=\delta^\mu{}_\nu$, buktikan
 $\pd_\lambda g^{\mu\nu}=-g^{\mu\alpha}g^{\nu\beta}\pd_\lambda g_{\alpha\beta}$.
 \item Turunkan simbol Christoffel seluruhnya dari $\nabla_\lambda g_{\mu\nu}=0$
 dan kondisi tanpa torsi.
 \item Hitung geodesik pada bidang polar datar dan tunjukkan bahwa hasilnya ekuivalen
 dengan garis lurus Kartesian.
 \item Buktikan bahwa $\dd(g_{\mu\nu}V^\mu W^\nu)/\dd\lambda=0$ untuk dua vektor
 yang ditransportasikan paralel.
 \item Ekspansi komutator $[\nabla_\mu,\nabla_\nu]V^\rho$ baris demi baris hingga
 seluruh turunan pertama $V$ saling hapus.
 \item Gunakan simetri Riemann untuk menghitung jumlah komponen independen di empat dimensi.
 \item Turunkan identitas Bianchi terkontraksi $\nabla_\mu G^{\mu\nu}=0$.
 \item Turunkan geodesic deviation dari keluarga geodesik dua-parameter.
 \item Variasikan $\sqrt{-g}$ menggunakan formula Jacobi untuk determinan.
 \item Turunkan EFE dari aksi Einstein--Hilbert, dengan menandai secara eksplisit
 suku batas.
 \item Turunkan persamaan Poisson dari komponen $00$ EFE linear.
 \item Hitung semua Christoffel independen dari ansatz Schwarzschild.
 \item Substitusikan Christoffel tersebut ke $R_{tt}$, $R_{rr}$, dan
 $R_{\theta\theta}$ tanpa perangkat simbolik.
 \item Dari persamaan vakum, peroleh $AB=1$, lalu $A=1-r_s/r$.
 \item Dari potensial efektif Schwarzschild, turunkan $r_{\rm ISCO}=6GM/c^2$.
 \item Inverskan blok $(0,\phi)$ metrik Kerr dan cek empat komponen identitas matriks.
 \item Tunjukkan langsung bahwa limit $a\to0$ metrik Kerr adalah Schwarzschild.
 \item Turunkan konservasi $\xi_\mu u^\mu$ dari persamaan Killing dan geodesik.
 \item Pisahkan persamaan Hamilton--Jacobi Kerr dan identifikasi konstanta Carter.
 \item Ekspansi $g_{0\phi}$ Kerr sampai orde pertama dalam $J$.
 \item Hitung curl dari $(\bm J\times\bm r)/r^3$ dan peroleh struktur medan dipol.
 \item Turunkan rata-rata orbit polar
 $\langle\hat r_i\hat r_j\rangle$ dan faktor $1/2$ pada presesi LT rata-rata.
 \item Buktikan secara analitik bahwa $C_u$, $C_S$, dan $C_\perp$ konstan sepanjang
 solusi kontinu sistem orbit--spin.
\end{enumerate}

\chapter*{Ringkasan Rantai yang Harus Dikuasai}
\addcontentsline{toc}{chapter}{Ringkasan Rantai yang Harus Dikuasai}

\begin{align*}
 \text{one-form}&:\quad \oneformp(V)=p_\mu V^\mu,
 \qquad \dual^\alpha(\e_\beta)=\delta^\alpha{}_\beta,\\
 \text{metrik}&:\quad V_\mu=g_{\mu\nu}V^\nu,\\
 \text{koneksi}&:\quad
 \Gamma^\rho{}_{\mu\nu}=\tfrac12g^{\rho\sigma}
 (g_{\sigma\mu,\nu}+g_{\sigma\nu,\mu}-g_{\mu\nu,\sigma}),\\
 \text{geodesik}&:\quad \ddot x^\rho+\Gamma^\rho{}_{\mu\nu}
 \dot x^\mu\dot x^\nu=0,\\
 \text{spin}&:\quad \dot S^\rho+\Gamma^\rho{}_{\mu\nu}u^\mu S^\nu=0,\\
 \text{Riemann}&:\quad [\nabla_\mu,\nabla_\nu]V^\rho
 =R^\rho{}_{\sigma\mu\nu}V^\sigma,\\
 \text{Ricci}&:\quad R_{\mu\nu}=R^\rho{}_{\mu\rho\nu},\\
 \text{Einstein}&:\quad G_{\mu\nu}=R_{\mu\nu}-\tfrac12Rg_{\mu\nu},\\
 \text{EFE}&:\quad G_{\mu\nu}=\frac{8\pi G}{c^4}T_{\mu\nu},\\
 \text{Schwarzschild}&:\quad A=1-\frac{2GM}{c^2r},\quad B=A^{-1},\\
 \text{Kerr}&:\quad a=\frac{J}{Mc},\quad g_{0\phi}\ne0,\\
 \text{geodetik}&:\quad
 \bm\Omega_{\rm geo}=\frac{3GM}{2c^2r^3}(\bm r\times\bm v),\\
 \text{Lense--Thirring}&:\quad
 \bm\Omega_{\rm LT}=\frac{G}{c^2r^3}
 [3\hat{\bm r}(\hat{\bm r}\cdot\bm J)-\bm J].
\end{align*}

\end{document}
